\documentclass[11pt,a4paper,sans]{moderncv}
\moderncvstyle{banking}
\moderncvcolor{blue}

% Force the name and section headings to render in moderncv blue (color1).
% Default banking leaves them black: moderncvstylebanking.sty's \colorlet
% copies (not aliases) the pre-scheme accent colour, so the name colours are
% frozen before \moderncvcolor runs. Re-let them after. \namefont is the hook
% every name-style macro routes through, so this also works on moderncv 2.3.1
% (Debian/Ubuntu apt), which has no \firstnamestyle/\lastnamestyle at all.
\renewcommand*{\namefont}{\fontsize{34}{36}\bfseries\upshape}
\colorlet{firstnamecolor}{color1}
\colorlet{lastnamecolor}{color1}
\colorlet{namecolor}{color1}
\renewcommand*{\sectionstyle}[1]{{\sectionfont\color{color1}#1}}

\usepackage[utf8]{inputenc}
% pdflatex fallback only (the documented engine is lualatex, which skips this
% branch). Without T1 font encoding pdflatex builds accented letters with
% \accent, and the PDF text layer stores them decomposed - `e` + U+0300 rather
% than U+00E8 - so an ATS keyword match on "Genève" fails while the page looks
% right. moderncv 2.5 loads T1 itself under pdflatex; 2.3.1 (Debian/Ubuntu apt)
% does not. \ifpdftex comes from iftex, which every moderncv version loads.
\ifpdftex\usepackage[T1]{fontenc}\fi
% moderncv loads hyperref itself in an \AtEndPreamble hook, so \hypersetup
% must go in an \AtEndPreamble of our own: on moderncv < 2.4 a top-level
% \usepackage{hyperref} clashes with the class's own
% \RequirePackage[unicode]{hyperref}. From 2.4.0 the class passes its options
% through \PassOptionsToPackage instead, which is what removes that clash.
\AtEndPreamble{\hypersetup{
    colorlinks=true,
    linkcolor=blue,
    filecolor=magenta,
    urlcolor=blue,
    pdftitle={[YOUR_NAME] - CV},
    % Keep pdfpagemode=UseNone: this block runs after moderncv's own
    % \AtEndPreamble (moderncv.cls sets pdfpagemode there), so a FullScreen
    % value here would win and open every CV in fullscreen presentation mode.
    pdfpagemode=UseNone,
}}
\usepackage[scale=0.80]{geometry}
\usepackage{import}

% Personal data
\name{[First]}{[Last]}
\address{[Your Address, City, Country]}{}{}
\phone[mobile]{[+XX XXXXXXXXXX]}
\email{[your.email@example.com]}
\extrainfo{\href{[LINKEDIN_URL]}{LinkedIn}, \href{[GITHUB_URL]}{GitHub}}

\begin{document}
\makecvtitle

\section{Profile}
[Concise elevator pitch summarizing your technical background, domain focus, and key strengths. 3-5 sentences highlighting what you bring to the role.]

\section{Core Competencies}
\begin{itemize}
\item \textbf{[Competency 1]}: [Brief description of skills, tools, and application]
\item \textbf{[Competency 2]}: [Brief description of skills, tools, and application]
\item \textbf{[Competency 3]}: [Brief description of skills, tools, and application]
\item \textbf{[Competency 4]}: [Brief description of skills, tools, and application]
\item \textbf{[Competency 5]}: [Brief description of skills, tools, and application]
\end{itemize}

\section{Professional Experience}
\cventry{[YYYY-Present]}{[Job Title]}{[Company]}{[City, Country]}{}{\vspace{1pt}
\begin{itemize}
\item {[Key responsibility or project focus area]}
\item {[Achievement 1: quantified result, metric, or business outcome]}
\item {[Achievement 2: technical milestone, architecture, or performance optimization]}
\item {[Achievement 3: cross-functional collaboration or leadership contribution]}
\end{itemize}
}
\vspace{3pt}
\cventry{[YYYY-YYYY]}{[Job Title]}{[Company]}{[City, Country]}{}{\vspace{1pt}
\begin{itemize}
\item {[Key responsibility or project focus area]}
\item {[Achievement 1: quantified result, metric, or business outcome]}
\item {[Achievement 2: technical milestone, architecture, or performance optimization]}
\item {[Achievement 3: tooling, automation, or efficiency improvement]}
\end{itemize}
}
\vspace{3pt}
\cventry{[YYYY-YYYY]}{[Job Title]}{[Company]}{[City, Country]}{}{\vspace{1pt}
\begin{itemize}
\item {[Key responsibility or project focus area]}
\item {[Achievement 1: quantified result, metric, or business outcome]}
\item {[Achievement 2: technical milestone, architecture, or performance optimization]}
\end{itemize}
}

\section{Education}
\cventry{[YYYY-YYYY]}{[Degree Level] in [Field of Study]}{[Institution]}{[City, Country]}{}{\vspace{1pt}
[Grade / Honors / Thesis topic if applicable. Relevant coursework and academic projects.]
}
\vspace{3pt}
\cventry{[YYYY-YYYY]}{[Degree Level] in [Field of Study]}{[Institution]}{[City, Country]}{}{\vspace{1pt}
[Grade / Honors / Thesis topic if applicable. Relevant coursework and academic projects.]
}

\section{Honors and Awards}
\begin{itemize}
\item \textbf{[Award or Recognition 1]} -- [Issuing Organization or Event, Year]
\item \textbf{[Award or Recognition 2]} -- [Issuing Organization or Event, Year]
\end{itemize}

\section{References}
\cventry{}{Available upon request.}{}{}{}{}

\end{document}
